\section{Density-one linear independence over square fields}
\label{sec:linear-independence}

We next show that \eqref{eq:LI} holds for all but \(o(1)\) of the parameters
in the simultaneous regime.  The square-field hypothesis is used at one
specific algebraic point: it makes the normalized Frobenius roots roots of a
rational polynomial.

For \(t\in U_g(\F_Q)\), write
\[
 P_{g,t}(T)=L(T,\chi_{g,t})=\prod_{j=1}^{2g}(1-\alpha_jT)
 \in\mathbb Z[T]
\]
and form its reciprocal polynomial
\begin{equation}
 R_{g,t}(X)=X^{2g}P_{g,t}(1/X)=\prod_{j=1}^{2g}(X-\alpha_j).
\label{eq:reciprocal-frob-polynomial}
\end{equation}
Suppose \(Q=s^2\), with \(s\in\mathbb Z_{>0}\).  The normalized inverse roots
\(\beta_j=\alpha_j/s\) are then the roots of
\begin{equation}
 \widetilde R_{g,t}(X)
 :=s^{-2g}R_{g,t}(sX)=\prod_{j=1}^{2g}(X-\beta_j)
 \in\mathbb Q[X].
\label{eq:normalized-rational-polynomial}
\end{equation}
Multiplication of every root by the nonzero rational number \(s^{-1}\)
changes neither the splitting field nor the permutation action of its Galois
group.  The roots of \eqref{eq:normalized-rational-polynomial} occur in
reciprocal pairs \(e^{i\vartheta_j},e^{-i\vartheta_j}\).
Write
\(W_{2g}=(\mathbb Z/2\mathbb Z)^g\rtimes S_g\) for the signed-permutation
Weyl group acting on these \(g\) reciprocal pairs.

\begin{lemma}[Maximal Weyl group implies angle--\(\pi\) independence]
\label{lem:maximal-galois-implies-LI}
Let \(g\geq2\), let \(Q\) be a square, and suppose that the splitting-field
Galois group of \(R_{g,t}\) is the full Weyl group \(W_{2g}\).  Then
\(\vartheta_1,\ldots,\vartheta_g,\pi\) are linearly independent over
\(\mathbb Q\).
\end{lemma}

\begin{proof}
Kowalski's relation-space theorem
\cite[Proposition~2.4(2), with \(m=1\)]{Kowalski2008} says that, under the
maximal \(W_{2g}\) hypothesis, the rational space of multiplicative relations
among the roots \(\beta_j\) is the reciprocal-pair space.  In particular, a
rational relation vector has equal coordinates on the two members of every
reciprocal pair.

Suppose that
\(\sum_{j=1}^g a_j\vartheta_j+b\pi=0\) with rational coefficients.  Clearing
denominators makes \(a_j,b\) integral.  Exponentiating twice removes the sign
coming from \(b\pi\) and gives
\begin{equation}
 \prod_{j=1}^g(e^{i\vartheta_j})^{2a_j}=1.
\label{eq:angle-multiplicative-relation}
\end{equation}
In the ordered reciprocal pairs
\((e^{i\vartheta_j},e^{-i\vartheta_j})\), the relation vector in
\eqref{eq:angle-multiplicative-relation} has coordinates \((2a_j,0)\).
Membership in the reciprocal-pair space forces \(2a_j=0\) for every \(j\).
The original additive relation then gives \(b=0\).  Notice that this argument
uses only the rational relation space; no saturation assertion about an
integral relation lattice is required.
\end{proof}

The required maximal-Galois estimate is already available for precisely our
one-parameter pencil.  We record its corrected form, since omitting the
erratum factor would make the genus dependence false.

\begin{proposition}[Exceptional parameters]
\label{prop:LI-exceptional-set}
There is an absolute constant \(A>0\) such that, whenever \(Q\) has odd
characteristic greater than \(2g+1\),
\begin{equation}
 \#\left\{t\in U_g(\F_Q):
   \operatorname{Gal}(R_{g,t}/\mathbb Q)\not\simeq W_{2g}\right\}
 \leq A g^2 Q^{1-\gamma_g}\log Q,
 \qquad
 \gamma_g=\frac1{4g^2+3g+5}.
\label{eq:kowalski-exceptional}
\end{equation}
Consequently, for square \(Q\),
\begin{equation}
 \Pr_{t\in U_g(\F_Q)}(t\text{ fails \eqref{eq:LI}})
 \ll g^2Q^{-\gamma_g}\log Q
\label{eq:LI-probability-bound}
\end{equation}
as soon as \(Q>4g\), with an absolute implied constant.
\end{proposition}

\begin{proof}
Kowalski's large-sieve theorem for
\(y^2=f(x)(x-t)\) gives \eqref{eq:kowalski-exceptional} with the stated
exponent and an absolute constant
\cite[Theorem~6.2]{Kowalski2006}.  Item~3 of the published erratum
\cite{KowalskiErratum} inserts the factor \(g^2\) in front of
\(Q^{1-\gamma_g}\log Q\).  Our polynomial \(f_g\) is squarefree because the
characteristic exceeds \(2g+1\), so the theorem applies.  Its statement uses
reducibility or splitting-field degree below \(2^gg!\).  The functional
equation already embeds the Galois action in the signed permutation group
\(W_{2g}\), whose order is \(2^gg!\); hence full degree is equivalent to
Galois group \(W_{2g}\).  Divide by
\(\lvert U_g(\F_Q)\rvert=Q-2g\) and invoke
\cref{lem:maximal-galois-implies-LI}.
\end{proof}

\begin{corollary}[Simultaneous density-one LI]
\label{cor:simultaneous-LI}
Let \(Q_g\) be odd square prime powers of characteristic \(p_g>2g+1\).  If
\begin{equation}
 \frac{\log Q_g}{g^2\log(g+2)}\longrightarrow\infty,
\label{eq:LI-growth-condition}
\end{equation}
then a uniformly chosen \(t_g\in U_g(\F_{Q_g})\) satisfies \eqref{eq:LI}
with probability \(1-o(1)\).
\end{corollary}

\begin{proof}
Taking logarithms in \eqref{eq:LI-probability-bound} gives at most
\begin{equation}
 O(\log g+\log\log Q_g)
 -\frac{\log Q_g}{4g^2+3g+5}.
\label{eq:LI-log-bound}
\end{equation}
Put \(x_g=(\log Q_g)/g^2\).  Hypothesis
\eqref{eq:LI-growth-condition} says \(x_g/\log(g+2)\to\infty\), while
\(\log\log Q_g=\log(g^2x_g)=o(x_g)\).  The negative term in
\eqref{eq:LI-log-bound} therefore dominates, and the expression tends to
\(-\infty\).
\end{proof}

Finally consider the explicit fields in \eqref{eq:explicit-fields}.  By
Bertrand's postulate,
\[
 2g+1<p_g<2(2g+1).
\]
The minimality of \(n_g\) gives
\begin{equation}
 g^2\log^2(g+2)\leq\log Q_g
 <g^2\log^2(g+2)+2\log p_g.
\label{eq:explicit-field-growth}
\end{equation}
Thus \(Q_g\) is an odd square, its characteristic exceeds \(2g+1\), and
\eqref{eq:LI-growth-condition} holds.  In particular, this completely
specified sequence has the density-one linear independence needed to replace
the actual endpoint law by \eqref{eq:finite-quenched-law}.  The much stronger
field growth is later used to absorb the heat-smoothing cost; the present
large-sieve step itself is only one component of that quantitative ledger.
